\chapter{Magnetic Fields and Electromagnetism}

There are actually two different magnetic fields:
the \textbf{\(\vecb{B}\) field} and the \textbf{\(\vecb{H}\) field}.
For this volume, we are only concerned with the \(\vecb{B}\) field.

\section{Ampère's Law}

Electric currents generate magnetic fields around themselves.
Ignoring edge effects, for a sufficiently long wire with constant current,
the magnetic field is given by \textbf{Ampere's law}:
\begin{equation}
    B = \frac{I \mu_0}{2 \pi r},
\end{equation}
where
\begin{itemize}
    \item \(B\) is the magnitude of the magnetic field.
    \item \(I\) is the current.
    \item \(\mu_0\) is the \textbf{permittivity of free space} or
        \textbf{vacuum permittivity}.
    \item \(r\) is the distance from the wire.
\end{itemize}

\section{Biot-Savart Law}

\textbf{Biot-Savart's law} is a generalization of Ampère's law
to wires that are not sufficiently long
and points where edge effects become non-negligible.
